\section*{Contoh}
\addcontentsline{toc}{section}{Contoh}

\begin{contoh}[Soal UTS dan jawaban]
\textbf{Soal.} Diberikan tata bahasa
\[
S \to aSb \mid \varepsilon.
\]
Tentukan tipe Chomsky, buat derivasi untuk \(aabb\), dan jelaskan
apakah tata bahasa ini LL(1).

\textbf{Jawaban.} Tata bahasa bertipe 2 (bebas konteks) karena aturan
berbentuk \(A \to \alpha\). Derivasi:
\[
S \Rightarrow aSb \Rightarrow aaSbb \Rightarrow aabb.
\]
Tata bahasa ini bukan LL(1) karena aturan \(S \to aSb\) dan
\(S \to \varepsilon\) memiliki FIRST yang tumpang tindih dengan
FOLLOW(\(S\)). Parser top-down memerlukan lookahead atau
transformasi.
\end{contoh}

\begin{contoh}[Soal UAS dan jawaban]
\textbf{Soal.} Diberikan program
\begin{lstlisting}[language=C]
int main() {
    int x = 3.5;
    break;
    return 0;
}
\end{lstlisting}
Sebutkan kesalahan semantik dan validasi, lalu tuliskan pesan yang
informatif.

\textbf{Jawaban.} Kesalahan semantik: \code{float} tidak dapat
ditetapkan ke \code{int} tanpa cast. Kesalahan validasi:
\code{break} di luar loop. Pesan:
\begin{lstlisting}[language=bash]
baris 1, kolom 9: tidak dapat menetapkan float ke int
baris 2, kolom 5: break hanya valid di dalam loop
\end{lstlisting}
\end{contoh}

\begin{contoh}[Pertanyaan demo proyek]
\begin{enumerate}[label=\arabic*.,leftmargin=*]
  \item Mengapa Anda memilih recursive descent atau Bison?
  \item Bagaimana parser menangani dangling else?
  \item Di mana tabel simbol menyimpan informasi scope?
  \item Bagaimana Anda menguji program tidak valid?
  \item Apa keterbatasan proyek Anda saat ini?
\end{enumerate}
Jawaban yang baik menunjukkan pemahaman trade-off, bukan sekadar
menyebut fitur.
\end{contoh}

\begin{contoh}[Kutipan esai refleksi]
``Pada awal semester, saya menganggap parsing hanya soal mencocokkan
token. Setelah mengerjakan proyek, saya memahami bahwa keputusan
desain grammar memengaruhi kualitas pesan kesalahan dan kemudahan
pengujian. Saya belajar memisahkan logika parsing dari logika
semantik dan menulis tes untuk setiap aturan.''
\end{contoh}

\begin{contoh}[Analisis ketercapaian]
Jika nilai kuis tinggi tetapi proyek rendah, kemungkinan mahasiswa
memahami konsep tetapi kesulitan integrasi. Jika proyek tinggi
tetapi ujian rendah, kemungkinan mahasiswa mahir praktik tetapi
lemah teori. Analisis ini membantu menentukan tindak lanjut yang
tepat.
\end{contoh}
